% compile: xelatex AUN3_7015906.tex
\documentclass{mko3}
\begin{document}
\mkothreetitle

% ----- หมวดที่ 1 ข้อมูลทั่วไป (รูปแบบ มคอ.3 AUN-QA) -----
\coursecode{7015906}
\coursename{ระเบียบวิธีวิจัยทางวิทยาศาสตร์และวิศวกรรมศาสตร์ (Science and Engineering Research Methodology)}
\program{หลักสูตรวิศวกรรมศาสตรมหาบัณฑิต สาขาวิชาวิศวกรรมคอมพิวเตอร์และปัญญาประดิษฐ์}
\coursecredits{3 (2-2-5)}
\coursegroup{วิชาเฉพาะด้านบังคับ}
\yearlevel{นักศึกษาปริญญาโท — ภาคต้น (ตามแผนการศึกษา)}
\semester{2 / 2568}
\instructor{ผศ.ดร.กาญจนา ดาวเด่น}
\contact{kanjana.dawden@uru.ac.th}
\coursedesc{ศึกษาทฤษฎีที่เกี่ยวข้องกับหลักการและระเบียบวิธีการวิจัย การทบทวนวรรณกรรม การวิเคราะห์ปัญหา การตั้งสมมุติฐานการวิจัย การแก้ปัญหาในงานวิจัย การเก็บรวบรวมข้อมูล การวิเคราะห์ข้อมูลโดยวิธีการทางสถิติ การทดสอบผลการวิเคราะห์และการสรุปผลการทดลอง การเขียนโครงการวิจัยและการเขียนรายงานวิจัย และจริยธรรมในการทำงานวิจัย}
\courseinfotable
\begin{clodomaintable}
  \cloitem{CLO1}{อธิบาย หลักการและระเบียบวิธีวิจัยทางวิทยาศาสตร์และวิศวกรรมศาสตร์ ได้อย่างถูกต้องตามหลักวิชาการ}{Cognitive}{2}
  \cloitem{CLO2}{วิเคราะห์ ปัญหาวิจัยและตั้งสมมติฐานการวิจัย ที่สอดคล้องกับบริบทของงานวิจัยทางวิศวกรรมคอมพิวเตอร์}{Cognitive}{4}
  \cloitem{CLO3}{จัดระบบ ข้อมูลจากการทบทวนวรรณกรรมและการเก็บรวบรวมข้อมูล เพื่อนำไปสู่การพัฒนาโจทย์วิจัย}{Affective}{4}
  \cloitem{CLO4}{แสดงออก ถึงการมีจริยธรรมในการทำวิจัยและการเขียนรายงาน ผ่านการอ้างอิงผลงานทางวิชาการอย่างถูกต้อง}{Affective}{5}
  \cloitem{CLO5}{ประยุกต์ใช้ วิธีการทางสถิติในการวิเคราะห์ข้อมูลและสรุปผลการวิจัย ได้อย่างเหมาะสมกับประเภทของข้อมูล}{Cognitive}{3}
\end{clodomaintable}

\begin{cloplotable}{5}{ & \textbf{PLO1} & \textbf{PLO2} & \textbf{PLO3} & \textbf{PLO4} & \textbf{PLO5}}
  \textbf{CLO1} &   & X &   &   &   \\ \hline
  \textbf{CLO2} &   &   & X &   &   \\ \hline
  \textbf{CLO3} &   &   & X &   &   \\ \hline
  \textbf{CLO4} &   &   &   & X &   \\ \hline
  \textbf{CLO5} &   &   &   &   & X \\ \hline
\end{cloplotable}

\begin{bloomtable}
  \bloomrow{CLO1}{2}{}{}{}
  \bloomrow{CLO2}{4}{}{}{}
  \bloomrow{CLO3}{}{4}{}{}
  \bloomrow{CLO4}{}{5}{}{}
  \bloomrow{CLO5}{3}{}{}{3}
\end{bloomtable}

\begin{contentanalysistable}
  \contentrow{หลักการและระเบียบวิธีวิจัยทางวิทยาศาสตร์และวิศวกรรมศาสตร์}{Cognitive 2}{อธิบาย}{ได้อย่างถูกต้องตามหลักวิชาการ (U) เพื่อนำไปใช้ในการทำวิจัย (CLO1) [PLO2]}{4}
  \contentrow{ปัญหาและสมมติฐานการวิจัย}{Cognitive 4}{วิเคราะห์}{ที่สอดคล้องกับบริบทของงานวิจัยทางวิศวกรรมคอมพิวเตอร์ (An) เพื่อกำหนดแนวทางการแก้ปัญหา (CLO2) [PLO3]}{4}
  \contentrow{ข้อมูลจากการทบทวนวรรณกรรมและการเก็บรวบรวมข้อมูล}{Affective 4}{จัดระบบ}{เพื่อนำไปสู่การพัฒนาโจทย์วิจัย (O) ที่มีคุณภาพและน่าเชื่อถือ (CLO3) [PLO3]}{4}
  \contentrow{จริยธรรมในการทำวิจัยและการเขียนรายงาน}{Affective 5}{แสดงออก}{ผ่านการอ้างอิงผลงานทางวิชาการอย่างถูกต้อง (In) ตามหลักจรรยาบรรณ (CLO4) [PLO4]}{4}
  \contentrow{วิธีการทางสถิติในการวิเคราะห์ข้อมูล}{Cognitive 3}{ประยุกต์ใช้}{และสรุปผลการวิจัยได้อย่างเหมาะสมกับประเภทของข้อมูล (Art) (CLO5) [PLO5]}{4}
\end{contentanalysistable}

\mkocoursedesc

\begin{courseactivities}
  \actitem{ศึกษาเอกสารและสื่อประกอบการสอน และทำแบบฝึกหัดเกี่ยวกับระเบียบวิธีวิจัย (CLO1)}
  \actitem{ค้นคว้าและวิเคราะห์บทความวิจัย และฝึกตั้งสมมุติฐานตามกรณีศึกษา (CLO2)}
  \actitem{ออกแบบและจัดระเบียบข้อมูลจากแหล่งต่างๆ และปฏิบัติการเก็บข้อมูล (CLO3)}
  \actitem{วิเคราะห์กรณีศึกษาทางจริยธรรม และฝึกเขียนอ้างอิงทางวิชาการ (CLO4)}
  \actitem{วิเคราะห์ข้อมูลจากชุดข้อมูลตัวอย่าง และฝึกใช้ซอฟต์แวร์สถิติ (CLO5)}
\end{courseactivities}

\begin{teachingtable}
  \teachrow{CLO1}{บรรยายเกี่ยวกับหลักการและแนวคิดของระเบียบวิธีวิจัย - อภิปรายกลุ่ม [IEL: ประสบการณ์เชิงบูรณาการ ลงมือปฏิบัติ พัฒนาท้องถิ่น]}{แบบทดสอบข้อเขียน - การตอบคำถามในชั้นเรียน}
  \teachrow{CLO2}{สาธิตการค้นคว้าข้อมูลวิจัย - ฝึกปฏิบัติการทบทวนวรรณกรรม [IEL: ประสบการณ์เชิงบูรณาการ ลงมือปฏิบัติ พัฒนาท้องถิ่น]}{รายงานการทบทวนวรรณกรรม - การนำเสนอแนวคิดงานวิจัย}
  \teachrow{CLO3}{แนะนำเทคนิคการจัดการข้อมูลวิจัย - ฝึกปฏิบัติการเก็บรวบรวมข้อมูล [IEL: ประสบการณ์เชิงบูรณาการ ลงมือปฏิบัติ พัฒนาท้องถิ่น]}{รายงานการจัดระบบข้อมูล - การนำเสนอ}
  \teachrow{CLO4}{บรรยายเกี่ยวกับจรรยาบรรณและจริยธรรมในการวิจัย - ยกตัวอย่างกรณีศึกษา [IEL: ประสบการณ์เชิงบูรณาการ ลงมือปฏิบัติ พัฒนาท้องถิ่น]}{รายงานกรณีศึกษาด้านจริยธรรม - การประเมินการเขียนอ้างอิง}
  \teachrow{CLO5}{สอนเทคนิคการใช้เครื่องมือทางสถิติ - ฝึกปฏิบัติการวิเคราะห์ข้อมูล [IEL: ประสบการณ์เชิงบูรณาการ ลงมือปฏิบัติ พัฒนาท้องถิ่น]}{รายงานผลการวิเคราะห์ข้อมูล - การนำเสนอผลการทดลอง}
\end{teachingtable}

\begin{assessmenttable}
  \assessrow{สอบก่อนเรียน (ถ้ามี)}{---}{0}
  \assessrow{สอบระหว่างเรียน}{CLO1, CLO2, CLO3}{30}
  \assessrow{สอบปลายภาค}{CLO1, CLO2, CLO3, CLO4, CLO5}{40}
  \assessrow{รายงาน / โครงงาน / กิจกรรม}{CLO3, CLO4, CLO5}{30}
\end{assessmenttable}

\begin{weeklyplan}
  \weekrow{1}{บทนำระเบียบวิธีวิจัย}{บรรยาย + ปฐมนิเทศ}{CLO1}{สอบก่อนเรียน}
  \weekrow{2}{ขั้นตอนการวิจัยทางวิทยาศาสตร์}{บรรยาย + อภิปราย}{CLO1}{แบบฝึกหัด}
  \weekrow{3}{การกำหนดปัญหาและวัตถุประสงค์}{บรรยาย + Workshop}{CLO2}{แบบฝึกหัด}
  \weekrow{4}{การทบทวนวรรณกรรม}{บรรยาย + ปฏิบัติการ}{CLO3}{รายงานย่อย}
  \weekrow{5}{การตั้งสมมติฐาน}{Workshop + อภิปราย}{CLO2}{สอบกลางภาค 1}
  \weekrow{6}{การออกแบบการวิจัย}{บรรยาย + Workshop}{CLO2}{แบบฝึกหัด}
  \weekrow{7}{การเก็บรวบรวมข้อมูล}{บรรยาย + กรณีศึกษา}{CLO3}{แบบฝึกหัด}
  \weekrow{8}{สถิติเชิงพรรณนา}{บรรยาย + Lab}{CLO5}{รายงานกลางภาค}
  \weekrow{9}{สถิติเชิงอนุมาน}{บรรยาย + Lab}{CLO5}{สอบกลางภาค 2}
  \weekrow{10}{จริยธรรมและการอ้างอิง}{อภิปรายกรณีศึกษา}{CLO4}{บทความสะท้อนคิด}
  \weekrow{11}{ออกแบบโครงการวิจัย}{PBL + ให้คำปรึกษา}{CLO2, CLO5}{ข้อเสนอโครงการ}
  \weekrow{12}{พัฒนาโครงการ (ช่วงที่ 1)}{Workshop}{CLO3, CLO5}{รายงานความก้าวหน้า}
  \weekrow{13}{พัฒนาโครงการ (ช่วงที่ 2)}{Workshop}{CLO5}{---}
  \weekrow{14}{สรุปและเตรียมนำเสนอ}{ฝึกนำเสนอ}{CLO2}{ร่างรายงาน}
  \weekrow{15}{นำเสนอโครงการและสอบปลายภาค}{นำเสนอ + สอบปลายภาค}{CLO1--CLO5}{สอบปลายภาค + โครงการ}
\end{weeklyplan}

% ----- หมวดที่ 5 Rubric -----
\begin{rubricsection}
\begin{subrubric}{CLO1}{อธิบาย หลักการและระเบียบวิธีวิจัยทางวิทยาศาสตร์และวิศวกรรมศาสตร์ ได้อย่างถูกต้องตามหลักวิชาการ}
  \rubricrow{4 (ดีเยี่ยม)}{อธิบายหลักการได้อย่างถูกต้องลึกซึ้ง เชื่อมโยงทฤษฎีกับบริบทวิชาชีพและวรรณกรรมที่เกี่ยวข้องได้อย่างเหมาะสม}
  \rubricrow{3 (ดี)}{อธิบายหลักการได้ถูกต้องส่วนใหญ่ ครอบคลุมประเด็นสำคัญตามมาตรฐานหลักสูตรปริญญาโท}
  \rubricrow{2 (พอใช้)}{อธิบายได้บางส่วน ยังขาดความลึกหรือการเชื่อมโยงเชิงวิชาการ}
  \rubricrow{1 (ต้องปรับปรุง)}{ไม่สามารถอธิบายได้อย่างถูกต้องหรือไม่ครบถ้วน}
\end{subrubric}

\begin{subrubric}{CLO2}{วิเคราะห์ ปัญหาวิจัยและตั้งสมมติฐานการวิจัย ที่สอดคล้องกับบริบทของงานวิจัยทางวิศวกรรมคอมพิวเตอร์}
  \rubricrow{4 (ดีเยี่ยม)}{วิเคราะห์เชิงวิพากษ์อย่างเป็นระบบ ใช้หลักฐาน/วรรณกรรมสนับสนุน สรุปและเสนอแนวทางได้ชัดเจน}
  \rubricrow{3 (ดี)}{วิเคราะห์ได้ถูกต้องส่วนใหญ่ มีเหตุผลประกอบเพียงพอตามมาตรฐานปริญญาโท}
  \rubricrow{2 (พอใช้)}{วิเคราะห์ได้ตื้น หรือยังขาดเหตุผล/หลักฐานสนับสนุน}
  \rubricrow{1 (ต้องปรับปรุง)}{วิเคราะห์ไม่ถูกต้องหรือไม่สามารถวิเคราะห์ได้}
\end{subrubric}

\begin{subrubric}{CLO3}{จัดระบบ ข้อมูลจากการทบทวนวรรณกรรมและการเก็บรวบรวมข้อมูล เพื่อนำไปสู่การพัฒนาโจทย์วิจัย}
  \rubricrow{4 (ดีเยี่ยม)}{จัดระบบแนวคิด/ความรู้ได้สมบูรณ์ เชื่อมโยงทฤษฎีและงานวิจัยได้อย่างน่าเชื่อถือ}
  \rubricrow{3 (ดี)}{จัดระบบได้ดี เชื่อมโยงแนวคิดได้ส่วนใหญ่}
  \rubricrow{2 (พอใช้)}{จัดระบบได้บางส่วน ยังไม่สอดคล้องหรือไม่ครบถ้วน}
  \rubricrow{1 (ต้องปรับปรุง)}{ไม่สามารถจัดระบบได้อย่างเหมาะสม}
\end{subrubric}

\begin{subrubric}{CLO4}{แสดงออก ถึงการมีจริยธรรมในการทำวิจัยและการเขียนรายงาน ผ่านการอ้างอิงผลงานทางวิชาการอย่างถูกต้อง}
  \rubricrow{4 (ดีเยี่ยม)}{แสดงทัศนคติ/จริยธรรมวิชาชีพได้ชัดเจน สม่ำเสมอ และสะท้อนคิดอย่างลึกซึ้ง}
  \rubricrow{3 (ดี)}{แสดงออกได้เหมาะสมส่วนใหญ่ มีการสะท้อนคิด}
  \rubricrow{2 (พอใช้)}{แสดงออกได้บางครั้ง ยังไม่ชัดเจนหรือไม่สม่ำเสมอ}
  \rubricrow{1 (ต้องปรับปรุง)}{ไม่แสดงออกหรือแสดงออกไม่เหมาะสม}
\end{subrubric}

\begin{subrubric}{CLO5}{ประยุกต์ใช้ วิธีการทางสถิติในการวิเคราะห์ข้อมูลและสรุปผลการวิจัย ได้อย่างเหมาะสมกับประเภทของข้อมูล}
  \rubricrow{4 (ดีเยี่ยม)}{ประยุกต์ใช้ได้อย่างถูกต้องและเหมาะสม แก้ปัญหาซับซ้อนได้ด้วยเหตุผลและหลักฐานสนับสนุน}
  \rubricrow{3 (ดี)}{ประยุกต์ใช้ได้ถูกต้องส่วนใหญ่ ตรงตามวัตถุประสงค์ของงาน}
  \rubricrow{2 (พอใช้)}{ประยุกต์ใช้ได้บางส่วน ยังไม่เหมาะสมหรือไม่สมบูรณ์}
  \rubricrow{1 (ต้องปรับปรุง)}{ไม่สามารถประยุกต์ใช้ได้หรือใช้ผิดแนวทาง}
\end{subrubric}
\end{rubricsection}


\begin{learningmaterials}
  \matitem{เอกสารประกอบการบรรยาย}
  \matitem{ตำรา: Creswell, \textit{Research Design}; Montgomery, \textit{Design and Analysis of Experiments}}
  \matitem{SPSS/R/Python สำหรับการวิเคราะห์สถิติ}
  \matitem{คู่มือจริยธรรมการวิจัยและการอ้างอิง}
  \matitem{ตัวอย่างรายงานวิจัยในสาขา CE/AI}
\end{learningmaterials}

\begin{supporttable}
  \supportrow{ห้องเรียน Smart Classroom}{1 ห้อง}{พร้อมใช้}{ทันสมัย}
  \supportrow{ห้องปฏิบัติการคอมพิวเตอร์}{1 ห้อง}{พร้อมใช้}{ทันสมัย}
  \supportrow{GPU Server / Cloud Computing Resources}{ครบถ้วน}{พร้อมใช้}{ทันสมัย}
  \supportrow{เอกสารประกอบการสอน, ตำรา, วารสารวิชาการ}{ครบถ้วน}{พร้อมใช้}{ทันสมัย}
\end{supporttable}

\signatureblock{ผศ.ดร.กาญจนา ดาวเด่น}{ผศ.ดร.พิทักษ์ คล้ายชม}

\end{document}